% 编译：xelatex demo-agent-skills.tex（跑两遍以稳定引用）
\documentclass{onepoem}   % 默认 xhs 开本；A4 版改为 \documentclass[a4]{onepoem}

\opheaderimg{../docs/one-poem-suffices/agent-skills/assets/Agent-Skills-Header.png}
\optitle{Agent Skills，一篇就够了。}
\opsubtitle{Composable Procedural Knowledge for Agents\\[3pt]Skills 本质是 Context，取用方式是 Just-in-Time.}
\opdate{2026.06}
\opread{50 MIN}

\begin{document}

\makecover

\section{What is a Skill?}

从格式标准上看，Skills 很简单：一个文件夹，一个 \code{SKILL.md}，再加可选的脚本和参考资料。它不像一个新的 Agent Framework，也不像 MCP 那样标准化外部能力的接入，更像一种可以被不同 Agent Harness 读取的轻量 packaging（参考 \weblink{https://agentskills.io/specification}{Agent Skills Specification}）。

\textbf{Skills 本质上就是 Context。}技术上你完全可以在 Skills 文件夹里放任何类型的上下文：既可以是指导性上下文，也可以是信息性上下文，事实数据、API 文档、甚至 RAG 索引。三层结构对应的 token 量级：

\begin{itemize}
  \item \textbf{第一层 metadata}（\code{name} + \code{description}）：约 100 token，启动时预加载进 System Prompt
  \item \textbf{第二层 \code{SKILL.md} 正文}：建议控制在 500 lines / 5,000 token 以内，被触发后才读
  \item \textbf{第三层 bundled files}（\code{references/}、\code{scripts/}）：无上限，按路径按需打开
\end{itemize}

\begin{tipbox}
我与一些朋友交流时发现，很多人不知道 Skills 的 metadata 会常驻 System Prompt，就像不知道 MCP 也会把 Tools 的定义注入上下文。如果忽略概念的包装，把两者都理解成 Context，那它们出现在上下文窗口里就是自然的事。
\end{tipbox}

一个最小的 Skill 长这样：

\begin{codeblock*}[my-skill/SKILL.md]{}
---
name: blog-review
description: Review blog drafts for style violations.
  Use when asked to "review my blog" or "check draft".
---

# Blog Review

1. Read the draft and the style guide in references/.
2. Flag banned patterns; suggest rewrites in place.
3. Output a prioritized list of findings.
\end{codeblock*}

\begin{questionbox}[scripts/ 是可执行的 tools，怎么能直接说 Skills 就是 Context？]
工具的定义和用法本身也是 Context，MCP 就是它的标准化协议，\code{scripts/} 可以看成它的本地版。Skill 从头到尾只负责提供 Context，``能跑''是 Harness/Sandbox 给的能力。
\end{questionbox}

\opfig[0.9]{../docs/one-poem-suffices/agent-skills/assets/skill-folder-context-semantics-gpt.png}{Skill 的三层结构：metadata 常驻，body 与 references 按需拉取}

\section{Why do we need Skills?}

Anthropic 对 Skill 的定义很简单：

\begin{blockquote}
Skills 是一组组织好的文件，为 Agent 打包可组合的程序性知识。换句话说，它们就是文件夹。这份简单是刻意设计的。\textit{``Skills are organized collections of files that package composable procedural knowledge for agents. In other words, they're folders. This simplicity is deliberate.''}
\end{blockquote}

\begin{importantbox}[Code is all we need]
我们过去以为不同领域的 Agent 会长得很不一样，每个都需要自己的工具和脚手架。但\textbf{底层的 Agent 其实比我们想象的更通用}：代码不只是一个用例，而是通往数字世界的通用接口。
\end{importantbox}

\subsection{Skill 和已有机制的区别}

\begin{optable}
\begin{tabularx}{\linewidth}{@{}ll>{\raggedright\arraybackslash}X@{}}
\toprule
\tabhead{机制} & \tabhead{加载时机} & \tabhead{典型场景} \\
\midrule
Prompt & 当前对话内输入 & 一次性请求、临场修正 \\
\code{CLAUDE.md} & 进入项目即全量加载 & 项目结构、测试命令、代码风格 \\
Skill & 任务语义命中才加载 & 发布检查、数据库 migration \\
\bottomrule
\end{tabularx}
\end{optable}

\opdivider

SkillsBench 的对照实验（84 个任务、11 个领域）给出了几项值得关注的结论：

\begin{optable}
\begin{tabularx}{\linewidth}{@{}l>{\raggedright\arraybackslash}X@{}}
\toprule
\tabhead{发现} & \tabhead{含义} \\
\midrule
人工精修的 Skill 平均提升 \textbf{16.2\%} & Skill 有用，但前提是由人来把关质量 \\
Healthcare 提升 ${\sim}52\%$，SWE 仅 4.5\% & 预训练覆盖越薄弱的领域，边际收益越大 \\
模型自生成的 Skill 提升 \textbf{0\%} & 自己编写再自己执行，没有正收益 \\
\bottomrule
\end{tabularx}
\end{optable}

\begin{notebox}[为什么叫 Procedural Knowledge？]
它来自认知科学，跟 declarative knowledge（``知道什么''）相对，指``知道怎么做''的那一类知识。不是每个 Skill 都是流程：一个只固定配色和字体的品牌设计 Skill，很难称为 Workflow Knowledge。
\end{notebox}

\begin{infobox}[可跳过的深水区]
接下来几段是基于``非官方''源码、对 Claude Code Memory 读写机制的拆解。如果你只关心差异结论，可以直接跳到下方，不影响后续阅读。就当是一次渐进式披露：这里是 reference file。
\end{infobox}

\begin{warnbox}
装第三方 Skill 约等于装可执行代码：来源不可信的 \code{scripts/} 存在供应链与注入风险，启用前先读一遍内容。
\end{warnbox}

回到文章开始的两个观点：\textbf{Skills 本质是 Context，Delivery 机制是 JIT / Pull}。程序性知识从来不缺（文档里有、人脑里有、训练数据里也有），缺的是让它按需到位的机制与规范。更多讨论见\weblink{https://keli-wen.github.io/One-Poem-Suffices/}{博客原文}。

\end{document}
